\documentclass[../../main.tex]{subfiles}

\begin{document}

\section{Riesz 表示定理}\label{section:ch19-riesz-representation}

令 $(X,\mathcal{M},\mu)$ 为 $\sigma$ 有限测度空间, $1\leqslant p<\infty$, 而 $q$ 是 $p$ 的共轭. 对 $g\in L^{q}(X,\mu)$, 定义泛函 $S_{g}:L^{p}(X,\mu)\to\mathbf{R}$ 为
\begin{align*}
\text{对所有 }f\in L^{p}(X,\mu),\quad S_{g}(f)=\int_{X}f\cdot g\,d\mu.
\end{align*}
根据 H\"older 不等式, $S_{g}$ 是 $L^{p}(X,\mu)$ 上的有界线性泛函且 $\|S_{g}\|=\|g\|_{q}$. Riesz 表示定理断言, $L^{p}(X,\mu)$ 上的每个有界线性泛函都具有这种形式.

\begin{lemma}\label{lemma:ch19-riesz-helper}
令 $(X,\mathcal{M},\mu)$ 为 $\sigma$ 有限测度空间而 $1\leqslant p<\infty$. 对于 $X$ 上的可积函数 $f$, 假定存在 $M\geqslant 0$ 使得对于 $X$ 上在有限测度集外消失的简单函数 $g$,
\begin{align}
\left|\int_{X}f\cdot g\,d\mu\right|\leqslant M\cdot\|g\|_{p}.\label{align:ch19-lemma}
\end{align}
则 $f$ 属于 $L^{q}(X,\mu)$, 其中 $q$ 是 $p$ 的共轭. 此外, $\|f\|_{q}\leqslant M$.
\end{lemma}

\begin{proof}
首先考虑 $p>1$ 的情形. 由于 $|f|$ 是非负可测函数而测度空间为 $\sigma$ 有限的, 根据简单逼近定理, 存在一个简单函数序列 $\{\varphi_{n}\}$, 该序列中的每个函数在有限测度集外消失, 在 $X$ 上逐点收敛于 $|f|$, 且对所有 $n$, $0\leqslant\varphi_{n}\leqslant|f|$. 由于在 $X$ 上 $\{\varphi_{n}^{q}\}$ 逐点收敛于 $|f|^{q}$, Fatou 引理告诉我们为证明 $|f|^{q}$ 是可积的且 $\|f\|_{q}\leqslant M$, 仅须证明对所有 $n$,
\begin{align}
\int_{X}\varphi_{n}^{q}\,d\mu\leqslant M^{q}.\label{align:ch19-lemma-proof1}
\end{align}
固定自然数 $n$. 为验证\eqref{align:ch19-lemma-proof1}, 我们估计泛函值 $\varphi_{n}^{q}$ 如下:
\begin{align}
\text{在 }X\text{ 上},\quad\varphi_{n}^{q}=\varphi_{n}\cdot\varphi_{n}^{q-1}\leqslant|f|\cdot\varphi_{n}^{q-1}=f\cdot\operatorname{sgn}(f)\cdot\varphi_{n}^{q-1}.\label{align:ch19-lemma-proof2}
\end{align}
定义简单函数 $g_{n}$ 为
\begin{align*}
\text{在 }X\text{ 上},\quad g_{n}=\operatorname{sgn}(f)\cdot\varphi_{n}^{q-1}.
\end{align*}
我们从\eqref{align:ch19-lemma-proof2}和\eqref{align:ch19-lemma}推出
\begin{align*}
\int_{X}\varphi_{n}^{q}\,d\mu\leqslant\int_{X}f\cdot g_{n}\,d\mu\leqslant M\cdot\|g_{n}\|_{p}.
\end{align*}
由于 $p$ 与 $q$ 共轭, $p(q-1)=q$, 因此,
\begin{align*}
\int_{X}|g_{n}|^{p}\,d\mu=\int_{X}\varphi_{n}^{p(q-1)}\,d\mu=\int_{X}\varphi_{n}^{q}\,d\mu.
\end{align*}
因此我们将上式重写为
\begin{align*}
\int_{X}\varphi_{n}^{q}\,d\mu\leqslant M\cdot\left[\int_{X}\varphi_{n}^{q}\,d\mu\right]^{1/p}.
\end{align*}
对每个 $n$, $\varphi_{n}^{q}$ 是在有限测度集外消失的简单函数, 因此它是可积的. 因此前一个积分不等式可重写为
\begin{align*}
\left[\int_{X}\varphi_{n}^{q}\,d\mu\right]^{1-1/p}\leqslant M.
\end{align*}
由于 $1-1/p=1/q$, 我们证明了\eqref{align:ch19-lemma-proof1}.

剩下来要考虑 $p=1$ 的情形. 我们必须证明 $M$ 是 $f$ 的本性上界. 用反证法. 若 $M$ 不是本性上界, 则存在某个 $\varepsilon>0$ 使得集合 $X_{\varepsilon}=\{x\in X\mid|f(x)|>M+\varepsilon\}$ 的测度非零. 由于 $X$ 是 $\sigma$ 有限的, 我们可以选取 $X_{\varepsilon}$ 的具有有限正测度的子集. 若令 $g$ 为这样的集合的特征函数, 得到与\eqref{align:ch19-lemma}矛盾的结论.
\end{proof}

\begin{theorem}[Riesz 表示定理]\label{theorem:ch19-riesz}
令 $(X,\mathcal{M},\mu)$ 为 $\sigma$ 有限测度空间, $1\leqslant p<\infty$, 而 $q$ 是 $p$ 的共轭. 对 $L^{p}(X,\mu)$ 上的每个有界线性泛函 $S$, 存在唯一的函数 $g\in L^{q}(X,\mu)$, 使得
\begin{align}
\text{对所有 }f\in L^{p}(X,\mu),\quad S(f)=\int_{X}f\cdot g\,d\mu.\label{align:ch19-riesz-theorem}
\end{align}
此外, $\|S\|=\|g\|_{q}$.
\end{theorem}

\begin{proof}
我们首先考虑 $\mu(X)<\infty$ 的情形. 令 $S:L^{p}(X,\mu)\to\mathbf{R}$ 为有界线性泛函. 定义可测集族 $\mathcal{M}$ 上的集函数 $\nu$ 为对 $E\in\mathcal{M}$,
\begin{align*}
\nu(E)=S(\chi_{E}).
\end{align*}
由于 $\chi_{E}\in L^{p}(X,\mu)$(因为 $\mu(X)<\infty$), 这是恰当定义的. 我们宣称 $\nu$ 是 $\mathcal{M}$ 上的有限测度. 有限可加性从 $S$ 的线性性得到. 对于可数可加性, 令 $\{E_{k}\}$ 为不交可测集族, 而 $E=\bigcup E_{k}$. 则 $\chi_{E}=\sum_{k=1}^{\infty}\chi_{E_{k}}$, 且在 $L^{p}$ 中 $\chi_{\bigcup_{k=1}^{n}E_{k}}\to\chi_{E}$. 因此
\begin{align*}
\nu(E)=\lim_{n\to\infty}S\left(\sum_{k=1}^{n}\chi_{E_{k}}\right)=\lim_{n\to\infty}\sum_{k=1}^{n}S(\chi_{E_{k}})=\sum_{k=1}^{\infty}\nu(E_{k}).
\end{align*}
此外, 若 $\mu(E)=0$, 则 $\chi_{E}=0$ a.e., 因此 $\nu(E)=S(0)=0$. 于是 $\nu\ll\mu$.

根据 Radon-Nikodym 定理, 存在 $g\in L^{1}(X,\mu)$, 使得对所有 $E\in\mathcal{M}$,
\begin{align*}
\nu(E)=\int_{E}g\,d\mu.
\end{align*}
因此, 对简单函数 $\psi=\sum c_{k}\chi_{E_{k}}$,
\begin{align*}
S(\psi)=\sum c_{k}\nu(E_{k})=\sum c_{k}\int_{E_{k}}g\,d\mu=\int_{X}\psi g\,d\mu.
\end{align*}
现在, 对任何简单函数 $\psi$,
\begin{align*}
\left|\int_{X}\psi g\,d\mu\right|=|S(\psi)|\leqslant\|S\|\cdot\|\psi\|_{p}.
\end{align*}
根据引理\cref{lemma:ch19-riesz-helper}, $g\in L^{q}(X,\mu)$ 且 $\|g\|_{q}\leqslant\|S\|$. 对于 $f\in L^{p}(X,\mu)$, 选取简单函数序列 $\{\psi_{n}\}$ 在 $L^{p}$ 中收敛于 $f$. 则
\begin{align*}
S(f)=\lim_{n\to\infty}S(\psi_{n})=\lim_{n\to\infty}\int_{X}\psi_{n}g\,d\mu=\int_{X}fg\,d\mu,
\end{align*}
其中最后一个等式从 H\"older 不等式得到. 此外, $\|S\|\leqslant\|g\|_{q}$, 因此 $\|S\|=\|g\|_{q}$.

对于 $\sigma$ 有限的情形, 将 $X$ 表示为有限测度集的不交并 $X=\bigcup X_{n}$, 对每个 $X_{n}$ 应用上述构造得到 $g_{n}$, 然后令 $g=\sum g_{n}\cdot\chi_{X_{n}}$. 唯一性从 H\"older 不等式得到.
\end{proof}

\begin{exercise}\label{exercise:ch19-06}
证明 Riesz 表示定理中函数 $g$ 的唯一性.
\end{exercise}

\begin{exercise}\label{exercise:ch19-07}
令 $(X,\mathcal{M},\mu)$ 为有限测度空间. 证明 $L^{\infty}(X,\mu)$ 是 $L^{1}(X,\mu)$ 的对偶空间的子空间, 但一般不是整个对偶空间.
\end{exercise}

\begin{exercise}\label{exercise:ch19-08}
令 $1<p<\infty$, 而 $q$ 是共轭指数. 证明 $L^{p}(X,\mu)$ 上的每个有界线性泛函 $S$ 满足 $\|S\|=\|g\|_{q}$, 其中 $g$ 是\eqref{align:ch19-riesz-theorem}中的代表元.
\end{exercise}

\end{document}